\documentclass{article}

\usepackage{amsthm, amsmath, amssymb}
\usepackage{listings}
\usepackage{xcolor}

% VS Code "Default Light+" inspired colors for Lean 4.
\definecolor{vscBg}{RGB}{255,255,255}        % white background
\definecolor{vscText}{RGB}{0,16,128}         % default identifier (dark blue)
\definecolor{vscKeyword}{RGB}{0,0,255}       % blue: def, theorem, axiom, ...
\definecolor{vscControl}{RGB}{175,0,219}     % purple: by, intro, exact, ...
\definecolor{vscComment}{RGB}{0,128,0}       % green: comments
\definecolor{vscString}{RGB}{163,21,21}      % dark red: strings
\definecolor{vscType}{RGB}{38,127,153}       % teal: Nat, Set, Prop, Finset, ...
\definecolor{vscFunction}{RGB}{121,94,38}    % brown: function names
\definecolor{vscNumber}{RGB}{9,134,88}       % green-blue: numerals
\definecolor{vscOperator}{RGB}{0,0,0}        % black: operators

\lstdefinelanguage{Lean4}{
  % Group 1: declaration keywords (blue)
  morekeywords=[1]{def,theorem,lemma,axiom,abbrev,instance,class,structure,
    inductive,namespace,end,open,import,attribute,deriving,where,variable,
    section},
  % Group 2: tactic and control keywords (purple)
  morekeywords=[2]{by,intro,intros,exact,apply,refine,rcases,obtain,simp,ext,
    sorry,rfl,decide,omega,ring,norm_num,dsimp,rw,set,constructor,rintro,
    interval_cases,simp_all,have,show,from,fun,let,in,if,then,else,match,
    with,this,calc,assumption,push_neg,by_cases,classical,split,refl,
    induction,cases},
  % Group 3: type names and important identifiers (teal)
  morekeywords=[3]{Nat,Int,Set,Prop,Type,Sort,Finset,Iff,And,Or,Not,Exists,
    Membership,Mem,Function,Subtype,Option,List,Array,Bool,True,False,
    Pythagorean,Prime,Infinite,IsHyp,Hd,Hp,Hu,H,PrimPythHyp,Finsupp},
  morecomment=[s][\color{vscComment}\itshape]{/-}{-/},
  morecomment=[l][\color{vscComment}\itshape]{--},
  morestring=[b][\color{vscString}]",
  basicstyle=\ttfamily\footnotesize\color{vscText},
  keywordstyle=[1]\color{vscKeyword},
  keywordstyle=[2]\color{vscControl},
  keywordstyle=[3]\color{vscType},
  identifierstyle=\color{vscText},
  showstringspaces=false,
  columns=fullflexible,
  keepspaces=true,
  literate=
    {ℕ}{{{\color{vscType}$\mathbb{N}$}}}1
    {ℤ}{{{\color{vscType}$\mathbb{Z}$}}}1
    {𝓗}{{$\mathcal{H}$}}1
    {∀}{{$\forall$}}1
    {∃}{{$\exists$}}1
    {∈}{{$\in$}}1
    {∉}{{$\notin$}}1
    {∧}{{$\wedge$}}1
    {∨}{{$\vee$}}1
    {¬}{{$\neg$}}1
    {→}{{$\to$}}1
    {←}{{$\leftarrow$}}1
    {↔}{{$\leftrightarrow$}}1
    {∏}{{$\prod$}}1
    {⇒}{{$\Rightarrow$}}1
    {≠}{{$\neq$}}1
    {≤}{{$\leq$}}1
    {≥}{{$\geq$}}1
    {₀}{{$_0$}}1
    {₁}{{$_1$}}1
    {₂}{{$_2$}}1
    {²}{{$^2$}}1
    {⟨}{{$\langle$}}1
    {⟩}{{$\rangle$}}1
    {↦}{{$\mapsto$}}1
    {·}{{$\cdot$}}1
    {×}{{$\times$}}1
    {ˢ}{{$^{\textrm{s}}$}}1
    {≡}{{$\equiv$}}1
    {∪}{{$\cup$}}1
    {∩}{{$\cap$}}1
    {⊆}{{$\subseteq$}}1
    {∅}{{$\emptyset$}}1
    {∣}{{$\mid$}}1
}
\usepackage[letterpaper, textwidth=6.5in, textheight=8.0in]{geometry}

\setlength{\parindent}{0.0in}
\setlength{\parskip}{1.5\baselineskip}

\title{A Conjecture on Primitive Pythagorean Triples}
\author{R. Scott McIntire}
\date{Aug 11, 2026}

\theoremstyle{plain}
\newtheorem{theorem}{Theorem}
\newtheorem{proposition}{Proposition}
\newtheorem{conjecture}{Conjecture}

\theoremstyle{definition}
\newtheorem{defn}{Definition}

\theoremstyle{remark}
\newtheorem{remark}{Remark}


\begin{document}
\maketitle


\section{Overview}
Some numbers carry a right triangle on their backs and some do not.
The number $5$ is the hypotenuse of the right triangle with sides
$(3,4,5)$; the numbers $7$ and $11$ are hypotenuses of no integer right
triangle at all. The number $25$ is a hypotenuse in exactly one
primitive way, $(7,24,25)$, while $65$ manages it twice:
$(16,63,65)$ and $(33,56,65)$. Two questions suggest themselves:
{\em which\/} numbers occur as hypotenuses of primitive Pythagorean
triples, and {\em in how many ways\/}? Classical number theory answers
both questions completely -- and the answers set up a third question,
about the primes among these hypotenuses, that appears to be genuinely
open.

In this note we discuss the set $\mathcal{H}$ of primitive Pythagorean hypotenuses --
the ``hypotenuse'' values that arise from primitive Pythagorean triples
(right triangles with integer sides sharing no common factor) --
recall the classical theorems describing its structure, and state
an open conjecture about the cumulative parity of its prime elements.

A related set, $\mathcal{H}_p \subset \mathcal{H}$, consists of those elements of
$\mathcal{H}$ that are prime -- the primitive Pythagorean hypotenuses that are
themselves prime numbers.

Two of the three observations below turn out to be classical theorems --
we recall them for orientation -- while the third remains an open conjecture
of mine, which is the focus of this note.
\begin{enumerate}
    \item \textbf{Partition of $\mathcal{H}$} (theorem): $\mathcal{H}$ is {\em partitioned\/}
      into two disjoint infinite sets: the set of all powers of Pythagorean primes,
      and the set of all hypotenuses arising from more than one primitive Pythagorean triple.
    \item \textbf{Infinitude of $\mathcal{H}_p$} (Dirichlet): The set of Pythagorean
      primes is infinite.
    \item \textbf{$\mathcal{H}_p \sim \mathbb{P}$} (conjecture): $\mathcal{H}_p$ is {\em similar\/}
      to and as {\em complex\/} as the primes when measured by certain
      cumulative-parity metrics.
\end{enumerate}


\section{Definitions}
After a few definitions we describe a conjecture regarding Primitive 
Pythagorean Triples. We use standard notation, associating the natural 
numbers with the symbol $\mathbb{N}$%
\footnote{We assume that $\mathbb{N}$ does {\em not\/} include $0$.}
and we use the symbol $\mathbb{P}$ to denote the set of prime numbers.


\begin{defn}
{\em Pythagorean Triples\/}, denoted $\mathcal{T}$, are triples of
numbers representing the lengths of the sides of right triangles where all the sides are integers.%
\footnote{Triples are taken as ordered, so $(3,4,5)$ and $(4,3,5)$ are
distinct elements of $\mathcal{T}$.}
Specifically,
\begin{equation}
  \mathcal{T} = \{ (x,y,z) \; | \; (x,y,z) \in \mathbb{N}^3 \; \wedge \; x^2 + y^2 = z^2 \}.
\end{equation}
\end{defn}

\begin{defn}
{\em Primitive Pythagorean Triples\/}, denoted $\mathcal{T}_p$, are
Pythagorean triples whose legs share no common factor:%
\footnote{Since any common divisor of $x$ and $y$ also divides $z$,
$\gcd(x,y) = 1$ is equivalent to $\gcd(x,y,z) = 1$, the usual primality condition.}
\begin{equation}
  \mathcal{T}_p = \{ (x,y,z) \; | \; (x,y,z) \in \mathbb{N}^3 \; \wedge \; \gcd(x,y) = 1 \; \wedge \; x^2 + y^2 = z^2 \}.
\end{equation}
\end{defn}

\begin{defn}
{\em Primitive Pythagorean Hypotenuses\/}, denoted $\mathcal{H}$, are defined by:
\begin{equation}
  \mathcal{H} = \{ z \; | \; \exists (x,y) \in \mathbb{N}^2, \;  (x,y,z) \in \mathcal{T}_p \}.
\end{equation}
This is the set of all possible hypotenuses of primitive Pythagorean triples.
\end{defn}

\begin{defn}
{\em Pythagorean Primes\/}, denoted $\mathcal{H}_p$, are defined by:
\begin{equation}
  \mathcal{H}_p = \{ z \; | \; \exists (x,y) \in \mathbb{N}^2, \, (x,y,z) \in \mathcal{T}_p \; \wedge \; z \in \mathbb{P} \}.
\end{equation}
This is the set of primes that occur as hypotenuses of primitive Pythagorean triples.
\end{defn}

\begin{defn}
{\em Duplicate Primitive Hypotenuses\/}, denoted $\mathcal{H}_d$, are defined by:
\begin{equation}
  \mathcal{H}_{d} = \{ z \; | \; \exists\, (x_1, x_2, y_1, y_2) \in \mathbb{N}^4, \;
      (x_1,y_1,z) \in \mathcal{T}_p \; \wedge \; (x_2, y_2, z) \in \mathcal{T}_p \; \wedge \; \{x_1, y_1\} \neq \{x_2, y_2\} \}.
\end{equation}
This is the set of primitive hypotenuses that arise from more than one
{\em distinct\/} primitive Pythagorean triple (i.e.\ not merely a swap of legs).
\end{defn}

\begin{defn}
{\em Primitive Hypotenuses Powers\/}, denoted $\mathcal{H}_u$, are defined by:
\begin{equation}
  \mathcal{H}_u = \{ z \; | \; \exists\,(p,n) \in \mathcal{H}_p \times \mathbb{N}, \; z = p^n \}.
\end{equation}
This is the set of all powers of Pythagorean primes.
\end{defn}


\textbf{Notes:}
\begin{itemize}
  \item The subscript $u$ in the name $\mathcal{H}_u$ stands for {\em unique\/}:
    by Theorem~\ref{thm:partition} below, the elements of $\mathcal{H}_u$ are
    exactly the hypotenuses arising from a unique primitive Pythagorean triple.
  \item $\mathcal{H}_p$ is a {\em proper\/} subset of the primes: $\mathcal{H}_p \subset \mathbb{P}$.
  \item If the sets above are treated as lists, they are indexed in sorted order.
\end{itemize}

The ground floor of this discussion is the parity structure of a primitive
triple: the hypotenuse is odd, one leg is even, and the even leg is
divisible by $4$. These facts follow from elementary parity arguments; an
alternative {\em geometric\/} proof -- via the inscribed circle of the right
triangle, which turns out to have integer radius -- appears in the companion
note~\cite{mcintire-prop}. The present paper starts where that one stops:
knowing the shape of a primitive triple, we ask which numbers occur as its
hypotenuse.

First we list a few facts and state a couple of known theorems before proceeding:
\begin{enumerate}
  \item The Pythagorean primes are exactly the primes of the form $4n + 1$.
      The key input is a theorem of Pierre de Fermat, known as
        Fermat's theorem on sums of two squares~\cite{hardy-wright}.
      According to this theorem, an odd prime number $p$ can be expressed as
        the sum of two squares if and only if
      $p$ is congruent to $1$ modulo $4$, i.e., $p = 4n + 1$ for some
        positive integer $n$.
      To pass from sums of two squares to hypotenuses, note that if
        $p = a^2 + b^2$ with $a > b$, then $p^2 = (a^2 - b^2)^2 + (2ab)^2$
        is a primitive representation, so $p \in \mathcal{H}_p$;
      conversely, a classical argument in the Gaussian integers
        $\mathbb{Z}[i]$ shows that every prime factor of a primitive
        hypotenuse is congruent to $1$ modulo $4$.
  \item Moreover, the distribution of primes of the form $4n + 1$ has
      been studied extensively.
      The German mathematician Peter Gustav Lejeune Dirichlet proved that
        there are infinitely many primes of the form $4n + 1$, a result known as
      Dirichlet's theorem on arithmetic progressions~\cite{davenport}.
      This theorem, which generalizes the case of Pythagorean primes, states
        that for any two positive coprime numbers $a$ and $d$,
      there are infinitely many primes of the form $a + nd$, where $n$ is a
        non-negative integer.
  \item In the case of Pythagorean primes, we can take $a = 1$ and $d = 4$ to
      get the form $4n + 1$.
        Dirichlet's theorem thus guarantees that there are infinitely
        many Pythagorean primes.
      This is a profound result, as it extends Euclid's theorem on the
      infinitude of the primes to arithmetic progressions.
\end{enumerate}
From these results we get the following characterization of Pythagorean primes:
\begin{theorem}\label{thm:hp-char}
$\mathcal{H}_p = \{ p \; | \; p \in \mathbb{P} \; \wedge \; p \equiv 1 \pmod{4}\}$.
\end{theorem}
And the following structure theorem for Pythagorean hypotenuses:
\begin{theorem}\label{thm:structure}
A natural number $z > 1$ belongs to $\mathcal{H}$ if and only if every prime
factor of $z$ lies in $\mathcal{H}_p$. Consequently, every $h \in \mathcal{H}$
factors uniquely over $\mathcal{H}_p$: there are a unique $N \in \mathbb{N}$,
unique $p_1 < p_2 < \cdots < p_N$ in $\mathcal{H}_p$, and unique
$(k_1, \ldots, k_N) \in \mathbb{N}^N$ such that
\begin{equation}
  h = \prod_{i=1}^N p_{i}^{k_i}.
\end{equation}
\end{theorem}

A classical consequence is the following partition theorem, which I had originally
stated as a conjecture but which follows from the count of primitive
representations of $n$ as a sum of two squares:
\begin{theorem}\label{thm:partition}
The set $\mathcal{H}$ is partitioned by $\mathcal{H}_d$ and $\mathcal{H}_u$:
\begin{enumerate}
  \item[(a)] $\hphantom{|\mathcal{H}_d|}\mathcal{H} = \mathcal{H}_{d} \cup \mathcal{H}_u$;
  \item[(b)] $\hphantom{|\mathcal{H}_d|}\emptyset = \mathcal{H}_{d} \cap \mathcal{H}_u$;
  \item[(c)] $\hphantom{\emptyset}|\mathcal{H}_d| = \infty$.
\end{enumerate}
\end{theorem}

\textbf{Proof sketch.} For $n \in \mathcal{H}$, the number of primitive Pythagorean
triples with hypotenuse $n$ equals $2^{k-1}$, where $k$ is the number of
distinct prime factors of $n$ (all of which lie in $\mathcal{H}_p$ by
Theorem~\ref{thm:structure}).
This count is a classical consequence of the unique factorization of $n$ in
the Gaussian integers~\cite{sierpinski}. Hence $n$ has exactly one primitive representation iff
$k = 1$, i.e.\ iff $n$ is a power of a single Pythagorean prime ($n \in \mathcal{H}_u$);
otherwise $n \in \mathcal{H}_d$. This gives the partition. The infinitude of
$\mathcal{H}_d$ follows from the infinitely many distinct products $p \cdot q$
with $p, q \in \mathcal{H}_p$ and $p \neq q$ (Dirichlet).\hfill$\square$

\begin{remark}
The set $\mathcal{H}_u$ is also infinite -- it contains $\mathcal{H}_p$, which
is infinite by Dirichlet's theorem -- so $\mathcal{H}$ is indeed partitioned
into two disjoint {\em infinite\/} sets, as described in the overview.
\end{remark}

It is worth naming the single strategic move behind all of the classical
results quoted above: {\em change the representation}. In the Gaussian
integers $\mathbb{Z}[i]$, the statement ``$z$ is a primitive Pythagorean
hypotenuse'' becomes ``$z^2 = (x + iy)(x - iy)$ with $x, y$ coprime'' -- a
statement about {\em factorization\/} -- and questions about triples become
bookkeeping about how $z$ factors: Fermat's theorem classifies which primes
survive the change of representation, and the $2^{k-1}$ count records how
many essentially distinct ways the prime factors of $z$ can be split
between $x + iy$ and its conjugate. Once the representation is right, the
theorems are inventory.

The following table shows the inventory for the smallest members of each
family, with the number of primitive triples doubling at each new distinct
prime factor (each row verified by direct computation):
\begin{center}
\begin{tabular}{c|c|l|c}
 distinct primes $k$ & triples $2^{k-1}$ & smallest example & lies in \\
\hline
 $1$ & $1$ & $5$: $\;(3,4,5)$ & $\mathcal{H}_u$ \\
 $1$ & $1$ & $25 = 5^2$: $\;(7,24,25)$ & $\mathcal{H}_u$ \\
 $2$ & $2$ & $65 = 5 \cdot 13$: $\;(16,63,65),\ (33,56,65)$ & $\mathcal{H}_d$ \\
 $3$ & $4$ & $1105 = 5 \cdot 13 \cdot 17$: $\;(47,1104,1105), \ldots$ & $\mathcal{H}_d$ \\
 $4$ & $8$ & $32045 = 5 \cdot 13 \cdot 17 \cdot 29$: $\;(716,32037,32045), \ldots$ & $\mathcal{H}_d$ \\
\end{tabular}
\end{center}
Theorem~\ref{thm:partition} is this table read down its last column: the
$k = 1$ rows -- the prime powers -- are exactly $\mathcal{H}_u$, and every
other row lands in $\mathcal{H}_d$.

\newpage
\section{McIntire's Pythagorean Triple Conjecture}

Having settled which numbers are hypotenuses and in how many ways, we turn
to the open question: how do the {\em prime\/} hypotenuses sit inside the
primes? Both $\mathcal{H}_p$ and $\mathbb{P}$ are infinite sorted lists,
and every prime $p \ge 5$ falls into one of two camps -- $p \equiv 1$ or
$p \equiv -1 \pmod 6$. Assign each camp a value:

\begin{defn}
The {\em Prime Parity\/} function, $\Pi$, is defined on any prime $p$ with $p \ge 5$ by:
\begin{equation}
  \Pi(p) = \begin{cases}
  \hphantom{-}1 & \exists n \in \mathbb{N}, \; p = 6n + 1 \\
            -1  & \exists n \in \mathbb{N}, \; p = 6n - 1
             \end{cases}
\end{equation}
\end{defn}

Now run a race between two walkers on the integers, both starting at $0$.
The first walker steps through the Pythagorean primes
$5, 13, 17, 29, 37, \ldots$ in order, moving up or down by
$\Pi$ at each step; the second steps through {\em all\/} primes
$\ge 5$: that is, $5, 7, 11, 13, 17, \ldots$.
After $m$ steps the walkers stand at the {\em cumulative parities\/}
$\sum_{i=1}^m \Pi(\mathcal{H}_p[i])$ and
$\sum_{i=3}^{m+2} \Pi(\mathbb{P}[i])$ respectively (the second sum starts
at index $3$ so as to skip the primes $2$ and $3$, on which $\Pi$ is
undefined; both sums thus range over primes $\ge 5$). The first few steps
of the race, with $D(m)$ the first walker's lead:
\begin{center}
\begin{tabular}{c|ccc|ccc|c}
 $m$ & $\mathcal{H}_p[m]$ & $\Pi$ & $\sum \Pi(\mathcal{H}_p[i])$
     & $\mathbb{P}[m+2]$ & $\Pi$ & $\sum \Pi(\mathbb{P}[i])$ & $D(m)$ \\
\hline
 $1$ & $5$  & $-1$ & $-1$ & $5$  & $-1$ & $-1$ & $\hphantom{-}0$ \\
 $2$ & $13$ & $+1$ & $\hphantom{-}0$ & $7$  & $+1$ & $\hphantom{-}0$ & $\hphantom{-}0$ \\
 $3$ & $17$ & $-1$ & $-1$ & $11$ & $-1$ & $-1$ & $\hphantom{-}0$ \\
 $4$ & $29$ & $-1$ & $-2$ & $13$ & $+1$ & $\hphantom{-}0$ & $-2$ \\
 $5$ & $37$ & $+1$ & $-1$ & $17$ & $-1$ & $-1$ & $\hphantom{-}0$ \\
 $6$ & $41$ & $-1$ & $-2$ & $19$ & $+1$ & $\hphantom{-}0$ & $-2$ \\
\end{tabular}
\end{center}
The natural question about any race is: how long must one wait for the
trailing runner to lead, and by how much? That waiting time is the
machinery this section needs:

\begin{defn}
The {\em First Deviation\/} function, $\Delta_p$, is defined by:%
\footnote{If the set on the right-hand side is empty, we set $\Delta_p(n) = \infty$.}
\begin{equation}
\Delta_p(n) = \min \left\{ m \;\middle|\; n \le \sum_{i=1}^m \Pi(\mathcal{H}_p[i]) - \sum_{i=3}^{m+2} \Pi(\mathbb{P}[i]) \right\}
\end{equation}
That is, $\Delta_p(n)$ is the number of steps until the Pythagorean walker
first leads by at least $n$.
\end{defn}

Two immediate observations. First, $\Delta_p$ is non-decreasing: the set
being minimized over shrinks as $n$ grows. Second, since both partial sums
are congruent to $m$ modulo $2$, the lead $D(m)$ is always {\em even\/};
consequently $\Delta_p(2k-1) = \Delta_p(2k)$ for all $k$. As an entry
fact: the Pythagorean walker first noses ahead only at step $41$ --
that is, $\Delta_p(1) = \Delta_p(2) = 41$.

We first record a fact that we had originally stated as a conjecture, but
which follows from the prime number theorem for arithmetic progressions:

\begin{proposition}\label{prop:mean-zero}
$\lim\limits_{n \rightarrow \infty} \dfrac{\sum\limits_{i=1}^n \Pi(\mathcal{H}_p[i])}{n} = 0$.
\end{proposition}

\textbf{Proof.} Every $p \in \mathcal{H}_p$ satisfies $p \equiv 1 \pmod 4$, and
$\Pi(p) = +1$ or $-1$ according as $p \equiv 1$ or $5 \pmod{12}$. By the prime
number theorem for arithmetic progressions~\cite{davenport}, the primes are
equidistributed among the $\varphi(12) = 4$ reduced residue classes
$1, 5, 7, 11 \pmod{12}$; in particular, of the first $n$ Pythagorean primes,
$n/2 + o(n)$ lie in each of the classes $1$ and $5 \pmod{12}$. The running sum
is the difference of these two counts, hence is $o(n)$.\hfill$\square$

Thus the $\pm 1$ values $\Pi(\mathcal{H}_p[i])$ average to zero, so
$\mathcal{H}_p$ is spread evenly between residues $1$ and $-1 \pmod 6$ in the
long run, mirroring $\mathbb{P}$ itself. Numerically, however, the cumulative
parity of $\mathcal{H}_p$ tends to run {\em lower\/} than that of $\mathbb{P}$:
over all primes below $2 \times 10^7$, the lead $D(m)$ is negative for
about $78\%$ of the steps and positive for only $21\%$.
The remaining open conjecture quantifies how often and how far the
$\mathcal{H}_p$ sum nonetheless overtakes the $\mathbb{P}$ sum:

\begin{conjecture}\label{conj:parity}
The cumulative parity of $\mathcal{H}_p$ exhibits the following non-trivial behavior:
\begin{enumerate}
  \item[(a)] $\forall n\in \mathbb{N}, \Delta_p(n) < \infty$;
  \item[(b)] $\lim\limits_{n \rightarrow \infty} \dfrac{\Delta_p(n)}{n^{\ln(2 \pi)}} = 41$.%
\footnote{The exponent $\ln(2\pi) \approx 1.838$ and the constant $41$ are
based on numerical fitting; we do not have a heuristic derivation, and the
reader should know how thin the evidence is. Computing with all primes below
$2 \times 10^7$ (the first $635{,}170$ Pythagorean primes, reaching leads
$n \le 136$), the ratio $\Delta_p(n)/n^{\ln(2\pi)}$ oscillates between
roughly $11$ and $112$ -- falling along each plateau of $\Delta_p$ and
jumping at each new record -- with no visible convergence over this range;
note also that $41$ is exactly $\Delta_p(1)$, so the constant is anchored
by the $n = 1$ data point. What the data defensibly supports is polynomial
growth with exponent near $1.8$; the precise limit and constant should be
treated as provisional, subject to revision as more data becomes
available.}
\end{enumerate}
\end{conjecture}

\begin{remark}
Part (b) implies part (a): $\Delta_p$ is non-decreasing, so if
$\Delta_p(n_0) = \infty$ for some $n_0$, then $\Delta_p(n) = \infty$ for all
$n \ge n_0$, contradicting the finiteness of the limit in (b). We state the
parts separately because (a) may well be provable without any understanding
of the growth rate in (b).
\end{remark}

Together with Proposition~\ref{prop:mean-zero}, the conjecture suggests that
the Pythagorean primes are ``as complex as'' the full set of primes when
measured by their parity behavior: deviations of any size $n$ do occur
($\Delta_p(n) < \infty$), and they grow at the polynomial rate
$n^{\ln(2\pi)}$. This ``mixing'' is what we mean by saying $\mathcal{H}_p$ is
non-trivial in a way similar to $\mathbb{P}$.

To return to where we began: the numbers $5, 13, 25, 65, \ldots$ that
carry primitive right triangles on their backs are completely classified --
which numbers occur (Theorem~\ref{thm:structure}), and in how many ways
(Theorem~\ref{thm:partition}). What remains open is subtler. The primes
among them, $5, 13, 17, 29, \ldots$, run their $\pm 1$ race against the
full sequence of primes, usually trailing; Conjecture~\ref{conj:parity}
asserts that they nonetheless take the lead by every margin eventually --
on a schedule that appears to be polynomial. Proving even part (a) would
be a satisfying close to the story that $(3,4,5)$ begins.

\newpage
\appendix
\section{A machine-checked Lean 4 proof of Theorem~\ref{thm:partition},
modulo three classical inputs}

The proof below has been mechanically checked by Lean 4
(toolchain \verb|leanprover/lean4:v4.29.1|) against Mathlib.
The full project lives at \verb|lean/PythHyp/| in this repository; the proof
file is\\\verb|lean/PythHyp/PythHyp/Basic.lean|.
The listing below is included directly from that source file, so it cannot
drift from what Lean actually checked.

The decomposition is honest: three classical results -- explicitly named
(F), (G-easy), and (G-uniq) -- are stated as \verb|axiom|s rather than
proved in this file, because their formalization, while certainly
possible in Mathlib, is a substantial project of its own. They are
classical theorems, not conjectures:
\begin{itemize}
  \item[(F)] {\em Fermat's characterization.\/} A positive integer $z > 1$
    is a primitive Pythagorean hypotenuse iff every prime factor of $z$ is
    $\equiv 1 \pmod 4$. The prime case is in Mathlib as
    \verb|Nat.Prime.sq_add_sq|; the general case follows from
    multiplicativity of the norm in $\mathbb Z[i]$.
  \item[(G-easy)] {\em Multiple factors yield duplicates.\/} If $z \in
    \mathcal H$ has at least two distinct prime factors, then $z$ admits
    at least two distinct unordered primitive representations
    $a^2 + b^2 = z^2$. This is the easy half of the classical
    $2^{k-1}$ count and follows from a single non-trivial
    factorization of $z$ in $\mathbb Z[i]$.
  \item[(G-uniq)] {\em Prime powers are unique.\/} If $z = p^n$ for a
    Pythagorean prime $p$ and $n \ge 1$, then $z$ admits only one
    unordered primitive representation. This is the harder half of the
    Gauss count.
\end{itemize}
Dirichlet's theorem on primes $\equiv 1 \pmod 4$ is {\em not\/} an axiom
in this proof: it is invoked directly as the Mathlib lemma
\verb|Nat.infinite_setOf_prime_modEq_one|.

Modulo the three classical inputs (F), (G-easy), (G-uniq), every step
of Theorem~\ref{thm:partition} -- parts (a), (b), and (c) -- is fully
formalized, with no \verb|sorry|. Lean reports a clean build, and
\verb|#print axioms PrimPythHyp.H_partition| confirms the accounting: the
final theorem depends on exactly the three axioms above, together with
Lean's standard \verb|propext|, \verb|Classical.choice|, and
\verb|Quot.sound|.

\lstinputlisting[language=Lean4, breaklines=true]{../lean/PythHyp/PythHyp/Basic.lean}

Once (F), (G-easy), (G-uniq) are themselves formalized, replacing the
\verb|axiom| declarations with \verb|theorem| proofs, this file would
become a self-contained, machine-checked proof of Theorem~\ref{thm:partition}
with no classical inputs assumed beyond what is already in Mathlib.

\begin{thebibliography}{9}

\bibitem{hardy-wright} G.~H. Hardy and E.~M. Wright,
  {\em An Introduction to the Theory of Numbers}, 6th ed.,
  Oxford University Press, 2008.
  (Fermat's two-square theorem, \S\S 20.1--20.3; Pythagorean triples, \S 13.2.)

\bibitem{davenport} H. Davenport,
  {\em Multiplicative Number Theory}, 3rd ed.,
  Graduate Texts in Mathematics 74, Springer, 2000.
  (Dirichlet's theorem and the prime number theorem for arithmetic progressions.)

\bibitem{sierpinski} W. Sierpi\'nski,
  {\em Pythagorean Triangles}, Dover Publications, 2003.
  (The $2^{k-1}$ count of primitive triples with a given hypotenuse.)

\bibitem{mcintire-prop} R.~S. McIntire,
  {\em An Alternative Proof of an Elementary Property of Pythagorean
  Triples}, companion note in this repository
  (\verb|tex/prop_pythag_triple.tex|), 2026.
  (The parity structure of a primitive triple via the inscribed circle.)

\end{thebibliography}

\end{document}


